\section{Représentation des programmes}
\subsection{Arbres de syntaxe abstraits}

\begin{frame}[fragile,t]
    \frametitle{Arbres de syntaxe abstraits (AST)}

    \begin{minted}{ocaml}
let rec sum (l : int list) : int =
    match l with
    | [] -> 0
    | x :: q -> x + sum q
    \end{minted}

    \pause

    \maxsizebox{\textwidth}{0.6\textheight}{
        \begin{forest}
            [{\texttt{FunctionDefinition(isRec: true, name: sum)}}
            [{body: \texttt{Match}},s sep=3cm
            [value : \texttt{Variable(l)}]
            [cases,s sep=3cm
            [\texttt{MatchCase}
            [{pattern : \texttt{EmptyList}},tier=forAlignment]
            [{expression : \texttt{Integer(0)}},tier=forAlignment]
            ]
            [\texttt{MatchCase},tier=forAlignment
            [pattern : \texttt{Concatenation}[...]]
            [expression : \texttt{Addition}
            [lhs : \texttt{Variable(x)}]
            [rhs : ...]
            ]
            ]
            ]
            ]
            [parameters: ...]
            [returnType: ...]
            ]
        \end{forest}
    }
\end{frame}

\begin{frame}[fragile,t]
    \frametitle{AST : Limitations}

    Les AST sont peu adaptés aux algorithmes demandant une connaissance de
    l'ordre d'évaluation du programme.

    \vspace{0.5cm}

    \begin{minipage}[t]{0.5\textwidth}
        \begin{minted}{ocaml}
a + b
        \end{minted}

        \begin{forest}
            [\texttt{Addition}
                [\texttt{Variable(a)}]
                [\texttt{Variable(b)}]]
        \end{forest}

        \only<2->{
            \begin{enumerate}
                \item \texttt{b}
                \item \texttt{a}
                \item \texttt{Addition}
            \end{enumerate}
        }
    \end{minipage}%
    \begin{minipage}[t]{0.5\textwidth}
        \begin{minted}{ocaml}
if a then b
        \end{minted}

        \begin{forest}
            [\texttt{If}
                [\texttt{Variable(a)}]
                [\texttt{Variable(b)}]]
        \end{forest}

        \only<2->{
            \begin{enumerate}
                \item \texttt{a}
                \item (Peut-être) \texttt{b}
            \end{enumerate}
        }
    \end{minipage}

    \vspace{0.5cm}

    \only<3->{
        \textcolor{orange}{$\blacktriangleright$} On préférerait une structure ne nécessitant pas de discriminer selon les étiquettes des nœuds.
    }
\end{frame}

\subsection{Formes linéaires}

\begin{frame}[fragile,t]
    \frametitle{Formes linéaires}

    \begin{itemize}
        \item La \important{forme linéaire} d'un programme est une liste
              d'étapes d'évaluation du programme.
        \item Chaque étape est désignée par un identifiant.

        \item Les étapes conditionnelles contiennent d'autres listes
              d'étapes.
    \end{itemize}
\end{frame}

\begin{frame}[fragile]
    \frametitle{Formes linéaires}

    \begin{minipage}{0.5\textwidth}
        \begin{minted}{ocaml}
if a = 2 then
    b * 3
else
    a - 2
        \end{minted}
    \end{minipage}%
    \begin{minipage}{0.5\textwidth}
        \small \begin{align*}
            a_1 :=\  & \texttt{2}                                      \\
            a_2 :=\  & \texttt{a}                                      \\
            a_3 :=\  & a_2 \texttt{ = } a_1                            \\
            a_4 :=\  & \texttt{if } a_3 \texttt{ then}                 \\
                     & \hspace{0.5cm} \begin{aligned}
                                          a_5 & := 3                    \\
                                          a_6 & := b                    \\
                                          a_7 & := a_6 \texttt{ * } a_5
                                      \end{aligned}    \\
                     & \texttt{else}                                   \\
                     & \hspace{0.5cm} \begin{aligned}
                                          a_8    & := 2                    \\
                                          a_9    & := a                    \\
                                          a_{10} & := a_9 \texttt{ - } a_8
                                      \end{aligned}
        \end{align*}
    \end{minipage}
\end{frame}

\subsection{Linéarisation et délinéarisation}

\begin{frame}[t]
    \frametitle{Linéarisation}

    On définit inductivement la linéarisation :

    \begin{itemize}
        \pause
        \item Linéariser(\texttt{a + b}) = Linéariser(\texttt{b}).Linéariser(\texttt{a}).[$a_k$ := $a_i$ \texttt{+} $a_j$] \\
              \pause
        \item Linéariser(\texttt{if a then b}) = \\
              \hspace{0.5cm} Linéariser(\texttt{a}).[$a_k$ := \texttt{if} $a_i$ \texttt{then} Linéariser(\texttt{b})]
              \pause
        \item ...
    \end{itemize}
    \small \textit{Avec $a_k$, $a_i$ et $a_j$ choisis en fonction des noms déjà utilisés}
\end{frame}

\begin{frame}[fragile,t]
    \frametitle{Délinéarisation}

    On peut aussi définir un algorithme transformant une forme linéaire en un
    programme \texttt{OCaml}:

    $$a_k := expr \longrightarrow \texttt{let a\_k = expr}$$

    \vspace{0.5cm}
    \pause

    \begin{minipage}{0.5\textwidth}
        Délinéariser\small$\left(\begin{aligned}
                    a_1 & := \texttt{a}                                \\
                    a_2 & := \texttt{2}                                \\
                    a_3 & := a_1 \texttt{ = } a_2                      \\
                    a_4 & := \texttt{if } a_3 \texttt{ then}           \\
                        & \hspace{0.5cm} \begin{aligned}
                                             a_5 & := 3                    \\
                                             a_6 & := b                    \\
                                             a_7 & := a_6 \texttt{ * } a_5
                                         \end{aligned} \\
                        & \texttt{else } ...
                \end{aligned}\right)$
    \end{minipage}=%
    \begin{minipage}{0.45\textwidth}
        \small \begin{minted}{ocaml}
let a_1 = a in
let a_2 = 2 in
let a_3 = a_1 = a_2 in
if a_3 then
    let a_5 = 3 in
    let a_6 = b in
    a_6 * a_5
else
    ...
            \end{minted}
    \end{minipage}
\end{frame}

\addtocounter{framenumber}{-1}
\begin{frame}[fragile,t]
    \frametitle{Délinéarisation}

    On peut aussi définir un algorithme transformant une forme linéaire en un
    programme \texttt{OCaml}:

    $$a_k := expr \longrightarrow \texttt{let a\_k = expr}$$

    \vspace{0.5cm}

    \begin{minipage}{0.5\textwidth}
        Délinéariser\small$\left(\begin{aligned}
                    a_1 & := \texttt{a}                                \\
                    a_2 & := \texttt{2}                                \\
                    a_3 & := a_1 \texttt{ = } a_2                      \\
                    a_4 & := \texttt{if } a_3 \texttt{ then}           \\
                        & \hspace{0.5cm} \begin{aligned}
                                             a_5 & := 3                    \\
                                             a_6 & := b                    \\
                                             a_7 & := a_6 \texttt{ * } a_5
                                         \end{aligned} \\
                        & \texttt{else } ...
                \end{aligned}\right)$
    \end{minipage}=%
    \begin{minipage}{0.45\textwidth}
        \small \begin{minted}{ocaml}
if a = 2 then
    b * 3
else
    a - 2
            \end{minted}
    \end{minipage}
\end{frame}

